%% ECON 282E -- Session 4, Deck B4: Fast and Accurate, Measured
%% Ported from QM Lec.8 F8-F14 (Howard), F19-F23 (monotonicity, concavity,
%% combining), F24-F36 (EGM), F48-F56 (errors, DH-M, grid sensitivity, standards).
%% THREE source problems fixed, all recorded in slides/SOURCES.md:
%%   (1) EGM's "no root-finding required" is false in the source's own
%%       Cobb-Douglas model; we use Carroll's cash-on-hand formulation, where it
%%       is true, and say so;
%%   (2) the "AFVRR (2006)" benchmark table is not from that paper (it has no EGM
%%       row and no time-iteration row) -- replaced by our own measured table;
%%   (3) the source claims 10x-100x on one frame and 50x-200x on another.
%% Howard's step count is written n_H: this course reserves H for the law of motion.
%% Numbers from tools/figures/s04_numbers.json, keys "rbc_quarterly",
%% "acceleration", "fig_euler_compare".

%% ECON 282E -- Foundations of Macroeconomics (UC Riverside, Fall 2026)
%% Shared Beamer preamble for every deck in slides/S01 ... slides/S10.
%%
%% Usage (a deck lives in slides/S0X/ and is compiled from that folder):
%%   %% ECON 282E -- Foundations of Macroeconomics (UC Riverside, Fall 2026)
%% Shared Beamer preamble for every deck in slides/S01 ... slides/S10.
%%
%% Usage (a deck lives in slides/S0X/ and is compiled from that folder):
%%   %% ECON 282E -- Foundations of Macroeconomics (UC Riverside, Fall 2026)
%% Shared Beamer preamble for every deck in slides/S01 ... slides/S10.
%%
%% Usage (a deck lives in slides/S0X/ and is compiled from that folder):
%%   \input{../preamble/course_preamble.tex}
%%   \decktitle{1}{A1}{Who I Am \& Why This Course}{September 24, 2026}
%%   \begin{document}
%%   \makedecktitle
%%   ...
%%
%% Adapted from Courses/AI-ECON-2026/source/shared_preamble.tex (Metropolis theme,
%% concept/caution/example/takeaway boxes, \sectiondivider). Changes for this course:
%%   * course title block (\decktitle / \makedecktitle) instead of \makebeamertitle
%%   * \zh{...} typesets Chinese (book titles) under pdflatex via CJKutf8
%%   * researchbox for the "Research in the wild" frames
%%   * section pages off; TikZ libraries and the course map loaded here
%% Build with pdflatex only (two passes); tools/build_slides.py does this for you.

\documentclass[10pt,english,aspectratio=169]{beamer}

\usepackage[T1]{fontenc}
\usepackage{textcomp}
\usepackage[utf8]{inputenc}
\usepackage{babel}
\usepackage{CJKutf8}

\usepackage{amsmath, amssymb, amsthm, graphicx}
\usepackage{booktabs, multirow, tabularx}
\usepackage{tikz}
\usetikzlibrary{positioning,arrows.meta,shapes,calc,fit,backgrounds}
\usepackage{pgfplots}
\pgfplotsset{compat=1.16}
\usepackage[most]{tcolorbox}
\tcbuselibrary{skins, breakable}
\usepackage{listings}
\usepackage{xcolor}

\lstset{
  basicstyle=\ttfamily\small,
  keywordstyle=\color{blue!80!black}\bfseries,
  commentstyle=\color{green!50!black}\itshape,
  stringstyle=\color{orange!80!black},
  backgroundcolor=\color{gray!8},
  frame=single,
  rulecolor=\color{gray!40},
  breaklines=true,
  showstringspaces=false,
  numbers=none,
}

\usetheme{metropolis}
\usefonttheme{professionalfonts}
\metroset{sectionpage=none}

\definecolor{LightGray}{RGB}{230,230,230}
\definecolor{RedTitle}{RGB}{0,0,200}
\definecolor{VeryLightGray}{RGB}{245,245,245}
\definecolor{DarkText}{RGB}{0,0,0}
\definecolor{UCRBlue}{RGB}{0,45,106}
\definecolor{UCRGold}{RGB}{241,171,0}

% Stage colours of the course arc (used by the course map and elsewhere)
\colorlet{StageTrunk}{blue!12}
\colorlet{StageBranch}{cyan!14}
\colorlet{StageMachine}{gray!18}
\colorlet{StageWall}{red!14}
\colorlet{StageTools}{orange!20}
\colorlet{StageData}{green!16}
\colorlet{StageAgents}{violet!14}

\hypersetup{bookmarks=false,breaklinks=true,pdfborder={0 0 0},colorlinks=true,
  linkcolor=DarkText,urlcolor=RedTitle}

\setbeamercolor{background canvas}{bg=white}
\setbeamercolor{title}{fg=RedTitle}
\setbeamercolor{frametitle}{bg=VeryLightGray,fg=DarkText}
\setbeamercolor{normal text}{fg=DarkText}
\setbeamercolor{alerted text}{fg=RedTitle}
\setbeamersize{text margin left=5mm, text margin right=5mm}
\addtobeamertemplate{frametitle}{}{\vspace{-0.5em}}

\setbeamertemplate{navigation symbols}{}
\setbeamertemplate{footline}{%
  \hspace{5mm}{\usebeamerfont{page number in head/foot}\color{black!45}%
  ECON 282E \,$\cdot$\, \insertshortdeck}\hfill%
  \usebeamercolor[fg]{page number in head/foot}%
  \usebeamerfont{page number in head/foot}%
  \insertframenumber\,/\,\inserttotalframenumber\kern1em\vskip2pt%
}

% ---------------------------------------------------------------------------
% Chinese text under pdflatex (book titles etc.)
\newcommand{\zh}[1]{\begin{CJK*}{UTF8}{gbsn}#1\end{CJK*}}

% ---------------------------------------------------------------------------
% Deck title block
%   \decktitle{session}{deck id}{title}{date}
\newcommand{\insertshortdeck}{}
\newcommand{\decksession}{}
\newcommand{\deckid}{}
\newcommand{\decktitletext}{}
\newcommand{\deckdate}{}
\newcommand{\decktitle}[4]{%
  \renewcommand{\decksession}{#1}%
  \renewcommand{\deckid}{#2}%
  \renewcommand{\decktitletext}{#3}%
  \renewcommand{\deckdate}{#4}%
  \renewcommand{\insertshortdeck}{Session #1 \,$\cdot$\, Deck #1.#2}%
  \title{#3}%
  \author{Zhigang Feng}%
  \date{#4}%
}
\newcommand{\makedecktitle}{%
  \begin{frame}[plain,noframenumbering]
    \begin{tikzpicture}[remember picture,overlay]
      \fill[UCRBlue] (current page.north west) rectangle ([yshift=-0.55cm]current page.north east);
      \fill[UCRGold] ([yshift=-0.55cm]current page.north west) rectangle ([yshift=-0.62cm]current page.north east);
      \node[anchor=west, text=white, font=\small\bfseries] at ([xshift=5mm,yshift=-0.275cm]current page.north west)
        {ECON 282E \,$\cdot$\, Foundations of Macroeconomics};
      \node[anchor=east, text=white, font=\small] at ([xshift=-5mm,yshift=-0.275cm]current page.north east)
        {UC Riverside \,$\cdot$\, Fall 2026};
    \end{tikzpicture}
    \vspace*{1.1cm}
    {\usebeamerfont{title}\color{black!55}\large Session \decksession\ \,$\cdot$\, Deck \decksession.\deckid\par}
    \vspace{0.25cm}
    {\usebeamerfont{title}\color{RedTitle}\LARGE\bfseries \decktitletext\par}
    \vspace{0.7cm}
    {\large Zhigang Feng\par}
    \vspace{0.15cm}
    {\color{black!65}\deckdate\par}
    \vfill
    {\footnotesize\itshape\color{black!55}Build it on paper $\;\to\;$ solve it on the machine $\;\to\;$ push it to the frontier with AI\par}
  \end{frame}}

% ---------------------------------------------------------------------------
% Section divider (plain frame, not counted)
\newcommand{\sectiondivider}[2]{%
\begin{frame}[plain,noframenumbering]
\begin{center}
\vfill
{\Large\textbf{#1}\par}
\vspace{0.35cm}
{\normalsize #2\par}
\vfill
\end{center}
\end{frame}
}

% ---------------------------------------------------------------------------
% Boxes
\newtcolorbox{conceptbox}[1][]{colback=blue!3, colframe=RedTitle!55, boxrule=0.35mm,
  arc=1mm, left=2mm, right=2mm, top=1mm, bottom=1mm, #1}
\newtcolorbox{cautionbox}[1][]{colback=red!3, colframe=red!50!black, boxrule=0.35mm,
  arc=1mm, left=2mm, right=2mm, top=1mm, bottom=1mm, #1}
\newtcolorbox{examplebox}[1][]{colback=green!3, colframe=green!45!black, boxrule=0.35mm,
  arc=1mm, left=2mm, right=2mm, top=1mm, bottom=1mm, #1}
\newtcolorbox{takeawaybox}[1][]{colback=VeryLightGray, colframe=RedTitle!55, boxrule=0.35mm,
  arc=1mm, left=2mm, right=2mm, top=1mm, bottom=1mm, #1}
\newtcolorbox{researchbox}[1][]{enhanced, colback=violet!4, colframe=violet!60!black,
  boxrule=0.35mm, arc=1mm, left=2mm, right=2mm, top=1mm, bottom=1mm,
  fonttitle=\bfseries\small, title={Research in the wild}, #1}

\newcommand{\compacttables}{\renewcommand{\arraystretch}{1.15}}
\newcommand{\src}[1]{{\tiny\color{black!55}#1}}

% ---------------------------------------------------------------------------
\input{../preamble/course_map.tex}

%%   \decktitle{1}{A1}{Who I Am \& Why This Course}{September 24, 2026}
%%   \begin{document}
%%   \makedecktitle
%%   ...
%%
%% Adapted from Courses/AI-ECON-2026/source/shared_preamble.tex (Metropolis theme,
%% concept/caution/example/takeaway boxes, \sectiondivider). Changes for this course:
%%   * course title block (\decktitle / \makedecktitle) instead of \makebeamertitle
%%   * \zh{...} typesets Chinese (book titles) under pdflatex via CJKutf8
%%   * researchbox for the "Research in the wild" frames
%%   * section pages off; TikZ libraries and the course map loaded here
%% Build with pdflatex only (two passes); tools/build_slides.py does this for you.

\documentclass[10pt,english,aspectratio=169]{beamer}

\usepackage[T1]{fontenc}
\usepackage{textcomp}
\usepackage[utf8]{inputenc}
\usepackage{babel}
\usepackage{CJKutf8}

\usepackage{amsmath, amssymb, amsthm, graphicx}
\usepackage{booktabs, multirow, tabularx}
\usepackage{tikz}
\usetikzlibrary{positioning,arrows.meta,shapes,calc,fit,backgrounds}
\usepackage{pgfplots}
\pgfplotsset{compat=1.16}
\usepackage[most]{tcolorbox}
\tcbuselibrary{skins, breakable}
\usepackage{listings}
\usepackage{xcolor}

\lstset{
  basicstyle=\ttfamily\small,
  keywordstyle=\color{blue!80!black}\bfseries,
  commentstyle=\color{green!50!black}\itshape,
  stringstyle=\color{orange!80!black},
  backgroundcolor=\color{gray!8},
  frame=single,
  rulecolor=\color{gray!40},
  breaklines=true,
  showstringspaces=false,
  numbers=none,
}

\usetheme{metropolis}
\usefonttheme{professionalfonts}
\metroset{sectionpage=none}

\definecolor{LightGray}{RGB}{230,230,230}
\definecolor{RedTitle}{RGB}{0,0,200}
\definecolor{VeryLightGray}{RGB}{245,245,245}
\definecolor{DarkText}{RGB}{0,0,0}
\definecolor{UCRBlue}{RGB}{0,45,106}
\definecolor{UCRGold}{RGB}{241,171,0}

% Stage colours of the course arc (used by the course map and elsewhere)
\colorlet{StageTrunk}{blue!12}
\colorlet{StageBranch}{cyan!14}
\colorlet{StageMachine}{gray!18}
\colorlet{StageWall}{red!14}
\colorlet{StageTools}{orange!20}
\colorlet{StageData}{green!16}
\colorlet{StageAgents}{violet!14}

\hypersetup{bookmarks=false,breaklinks=true,pdfborder={0 0 0},colorlinks=true,
  linkcolor=DarkText,urlcolor=RedTitle}

\setbeamercolor{background canvas}{bg=white}
\setbeamercolor{title}{fg=RedTitle}
\setbeamercolor{frametitle}{bg=VeryLightGray,fg=DarkText}
\setbeamercolor{normal text}{fg=DarkText}
\setbeamercolor{alerted text}{fg=RedTitle}
\setbeamersize{text margin left=5mm, text margin right=5mm}
\addtobeamertemplate{frametitle}{}{\vspace{-0.5em}}

\setbeamertemplate{navigation symbols}{}
\setbeamertemplate{footline}{%
  \hspace{5mm}{\usebeamerfont{page number in head/foot}\color{black!45}%
  ECON 282E \,$\cdot$\, \insertshortdeck}\hfill%
  \usebeamercolor[fg]{page number in head/foot}%
  \usebeamerfont{page number in head/foot}%
  \insertframenumber\,/\,\inserttotalframenumber\kern1em\vskip2pt%
}

% ---------------------------------------------------------------------------
% Chinese text under pdflatex (book titles etc.)
\newcommand{\zh}[1]{\begin{CJK*}{UTF8}{gbsn}#1\end{CJK*}}

% ---------------------------------------------------------------------------
% Deck title block
%   \decktitle{session}{deck id}{title}{date}
\newcommand{\insertshortdeck}{}
\newcommand{\decksession}{}
\newcommand{\deckid}{}
\newcommand{\decktitletext}{}
\newcommand{\deckdate}{}
\newcommand{\decktitle}[4]{%
  \renewcommand{\decksession}{#1}%
  \renewcommand{\deckid}{#2}%
  \renewcommand{\decktitletext}{#3}%
  \renewcommand{\deckdate}{#4}%
  \renewcommand{\insertshortdeck}{Session #1 \,$\cdot$\, Deck #1.#2}%
  \title{#3}%
  \author{Zhigang Feng}%
  \date{#4}%
}
\newcommand{\makedecktitle}{%
  \begin{frame}[plain,noframenumbering]
    \begin{tikzpicture}[remember picture,overlay]
      \fill[UCRBlue] (current page.north west) rectangle ([yshift=-0.55cm]current page.north east);
      \fill[UCRGold] ([yshift=-0.55cm]current page.north west) rectangle ([yshift=-0.62cm]current page.north east);
      \node[anchor=west, text=white, font=\small\bfseries] at ([xshift=5mm,yshift=-0.275cm]current page.north west)
        {ECON 282E \,$\cdot$\, Foundations of Macroeconomics};
      \node[anchor=east, text=white, font=\small] at ([xshift=-5mm,yshift=-0.275cm]current page.north east)
        {UC Riverside \,$\cdot$\, Fall 2026};
    \end{tikzpicture}
    \vspace*{1.1cm}
    {\usebeamerfont{title}\color{black!55}\large Session \decksession\ \,$\cdot$\, Deck \decksession.\deckid\par}
    \vspace{0.25cm}
    {\usebeamerfont{title}\color{RedTitle}\LARGE\bfseries \decktitletext\par}
    \vspace{0.7cm}
    {\large Zhigang Feng\par}
    \vspace{0.15cm}
    {\color{black!65}\deckdate\par}
    \vfill
    {\footnotesize\itshape\color{black!55}Build it on paper $\;\to\;$ solve it on the machine $\;\to\;$ push it to the frontier with AI\par}
  \end{frame}}

% ---------------------------------------------------------------------------
% Section divider (plain frame, not counted)
\newcommand{\sectiondivider}[2]{%
\begin{frame}[plain,noframenumbering]
\begin{center}
\vfill
{\Large\textbf{#1}\par}
\vspace{0.35cm}
{\normalsize #2\par}
\vfill
\end{center}
\end{frame}
}

% ---------------------------------------------------------------------------
% Boxes
\newtcolorbox{conceptbox}[1][]{colback=blue!3, colframe=RedTitle!55, boxrule=0.35mm,
  arc=1mm, left=2mm, right=2mm, top=1mm, bottom=1mm, #1}
\newtcolorbox{cautionbox}[1][]{colback=red!3, colframe=red!50!black, boxrule=0.35mm,
  arc=1mm, left=2mm, right=2mm, top=1mm, bottom=1mm, #1}
\newtcolorbox{examplebox}[1][]{colback=green!3, colframe=green!45!black, boxrule=0.35mm,
  arc=1mm, left=2mm, right=2mm, top=1mm, bottom=1mm, #1}
\newtcolorbox{takeawaybox}[1][]{colback=VeryLightGray, colframe=RedTitle!55, boxrule=0.35mm,
  arc=1mm, left=2mm, right=2mm, top=1mm, bottom=1mm, #1}
\newtcolorbox{researchbox}[1][]{enhanced, colback=violet!4, colframe=violet!60!black,
  boxrule=0.35mm, arc=1mm, left=2mm, right=2mm, top=1mm, bottom=1mm,
  fonttitle=\bfseries\small, title={Research in the wild}, #1}

\newcommand{\compacttables}{\renewcommand{\arraystretch}{1.15}}
\newcommand{\src}[1]{{\tiny\color{black!55}#1}}

% ---------------------------------------------------------------------------
%% Course map: the recurring "you are here" slide for ECON 282E.
%% Loaded by course_preamble.tex. Pattern follows AI-ECON-2026 lec01_map.tex, but the map
%% shows the whole ten-session course rather than one lecture.
%%
%%   \CourseMap{s}{label}   s = 1..10 highlights session s and prints "you are here: label"
%%                          (e.g. \CourseMap{1}{1.A2}); earlier sessions get a check mark.
%%   \CourseMap{0}{}        full overview, nothing highlighted.
%%
%% Note: TikZ option lists are comma-split before TeX conditionals run, so every \ifnum
%% below sits outside [...] option lists.

\newcommand{\CourseMap}[2]{%
\begin{center}
\begin{tikzpicture}[
    sess/.style={rectangle, rounded corners=2pt, draw=black!45, minimum width=1.34cm,
        minimum height=1.12cm, text width=1.26cm, align=center, inner sep=1pt},
    stage/.style={font=\tiny\bfseries, text=black!70},
]
  \foreach \i/\d/\t/\c in {%
      1/Sep 24/Intro \\ Growth/StageTrunk,
      2/Oct 1/HA $\cdot$ OLG \\ Policy/StageBranch,
      3/Oct 8/Num.\ DP \\ Python/StageMachine,
      4/Oct 15/Numerics \\ I/StageMachine,
      5/Oct 22/Numerics \\ II/StageMachine,
      6/Oct 29/The wall \\ \textbf{Midterm}/StageWall,
      7/Nov 5/ML \\ Deep nets/StageTools,
      8/Nov 12/RL \\ HA $+$ DL/StageTools,
      9/Nov 19/Text \\ LLMs/StageData,
      10/Dec 3/Agents \\ \textbf{Projects}/StageAgents} {
    \node[sess, fill=\c] (n\i) at ({1.46*(\i-1)}, 0)
      {{\scriptsize\bfseries S\i}\\[-1pt]{\tiny\color{black!60}\d}\\[1pt]{\tiny \t}};
    \ifnum\i<#1
      \node[font=\scriptsize, text=green!45!black] at ([xshift=-2.5mm,yshift=-2.2mm]n\i.north east) {$\checkmark$};
    \fi
    \ifnum\i=#1
      \draw[RedTitle, very thick, rounded corners=2pt]
        ([xshift=-0.6mm,yshift=-0.6mm]n\i.south west) rectangle ([xshift=0.6mm,yshift=0.6mm]n\i.north east);
      % keep the label inside the picture at both ends of the row
      \ifnum\i<3
        \node[font=\scriptsize\bfseries, text=RedTitle, anchor=south west] at ([yshift=1.2mm]n\i.north west)
          {$\blacktriangledown$\,you are here: #2};
      \else\ifnum\i>8
        \node[font=\scriptsize\bfseries, text=RedTitle, anchor=south east] at ([yshift=1.2mm]n\i.north east)
          {you are here: #2\,$\blacktriangledown$};
      \else
        \node[font=\scriptsize\bfseries, text=RedTitle, anchor=south] at ([yshift=1.2mm]n\i.north)
          {you are here: #2\,$\blacktriangledown$};
      \fi\fi
    \fi
  }
  % stage band under the sessions
  \foreach \a/\b/\lab/\c in {1/1/trunk/StageTrunk, 2/2/branches/StageBranch,
      3/5/the machine $\cdot$ classical methods/StageMachine, 6/6/the wall/StageWall,
      7/8/new tools/StageTools, 9/9/new data/StageData, 10/10/colleagues/StageAgents} {
    \fill[\c] ([yshift=-1.5mm]n\a.south west) rectangle ([yshift=-4.6mm]n\b.south east);
    \path ([yshift=-1.5mm]n\a.south west) -- ([yshift=-4.6mm]n\b.south east)
      node[midway, stage] {\lab};
  }
  \node[font=\footnotesize\itshape, text=black!70] at ([yshift=-9.5mm]$(n1.south)!0.5!(n10.south)$)
    {Build it on paper $\;\to\;$ solve it on the machine $\;\to\;$ push it to the frontier with AI};
\end{tikzpicture}
\end{center}
}


%%   \decktitle{1}{A1}{Who I Am \& Why This Course}{September 24, 2026}
%%   \begin{document}
%%   \makedecktitle
%%   ...
%%
%% Adapted from Courses/AI-ECON-2026/source/shared_preamble.tex (Metropolis theme,
%% concept/caution/example/takeaway boxes, \sectiondivider). Changes for this course:
%%   * course title block (\decktitle / \makedecktitle) instead of \makebeamertitle
%%   * \zh{...} typesets Chinese (book titles) under pdflatex via CJKutf8
%%   * researchbox for the "Research in the wild" frames
%%   * section pages off; TikZ libraries and the course map loaded here
%% Build with pdflatex only (two passes); tools/build_slides.py does this for you.

\documentclass[10pt,english,aspectratio=169]{beamer}

\usepackage[T1]{fontenc}
\usepackage{textcomp}
\usepackage[utf8]{inputenc}
\usepackage{babel}
\usepackage{CJKutf8}

\usepackage{amsmath, amssymb, amsthm, graphicx}
\usepackage{booktabs, multirow, tabularx}
\usepackage{tikz}
\usetikzlibrary{positioning,arrows.meta,shapes,calc,fit,backgrounds}
\usepackage{pgfplots}
\pgfplotsset{compat=1.16}
\usepackage[most]{tcolorbox}
\tcbuselibrary{skins, breakable}
\usepackage{listings}
\usepackage{xcolor}

\lstset{
  basicstyle=\ttfamily\small,
  keywordstyle=\color{blue!80!black}\bfseries,
  commentstyle=\color{green!50!black}\itshape,
  stringstyle=\color{orange!80!black},
  backgroundcolor=\color{gray!8},
  frame=single,
  rulecolor=\color{gray!40},
  breaklines=true,
  showstringspaces=false,
  numbers=none,
}

\usetheme{metropolis}
\usefonttheme{professionalfonts}
\metroset{sectionpage=none}

\definecolor{LightGray}{RGB}{230,230,230}
\definecolor{RedTitle}{RGB}{0,0,200}
\definecolor{VeryLightGray}{RGB}{245,245,245}
\definecolor{DarkText}{RGB}{0,0,0}
\definecolor{UCRBlue}{RGB}{0,45,106}
\definecolor{UCRGold}{RGB}{241,171,0}

% Stage colours of the course arc (used by the course map and elsewhere)
\colorlet{StageTrunk}{blue!12}
\colorlet{StageBranch}{cyan!14}
\colorlet{StageMachine}{gray!18}
\colorlet{StageWall}{red!14}
\colorlet{StageTools}{orange!20}
\colorlet{StageData}{green!16}
\colorlet{StageAgents}{violet!14}

\hypersetup{bookmarks=false,breaklinks=true,pdfborder={0 0 0},colorlinks=true,
  linkcolor=DarkText,urlcolor=RedTitle}

\setbeamercolor{background canvas}{bg=white}
\setbeamercolor{title}{fg=RedTitle}
\setbeamercolor{frametitle}{bg=VeryLightGray,fg=DarkText}
\setbeamercolor{normal text}{fg=DarkText}
\setbeamercolor{alerted text}{fg=RedTitle}
\setbeamersize{text margin left=5mm, text margin right=5mm}
\addtobeamertemplate{frametitle}{}{\vspace{-0.5em}}

\setbeamertemplate{navigation symbols}{}
\setbeamertemplate{footline}{%
  \hspace{5mm}{\usebeamerfont{page number in head/foot}\color{black!45}%
  ECON 282E \,$\cdot$\, \insertshortdeck}\hfill%
  \usebeamercolor[fg]{page number in head/foot}%
  \usebeamerfont{page number in head/foot}%
  \insertframenumber\,/\,\inserttotalframenumber\kern1em\vskip2pt%
}

% ---------------------------------------------------------------------------
% Chinese text under pdflatex (book titles etc.)
\newcommand{\zh}[1]{\begin{CJK*}{UTF8}{gbsn}#1\end{CJK*}}

% ---------------------------------------------------------------------------
% Deck title block
%   \decktitle{session}{deck id}{title}{date}
\newcommand{\insertshortdeck}{}
\newcommand{\decksession}{}
\newcommand{\deckid}{}
\newcommand{\decktitletext}{}
\newcommand{\deckdate}{}
\newcommand{\decktitle}[4]{%
  \renewcommand{\decksession}{#1}%
  \renewcommand{\deckid}{#2}%
  \renewcommand{\decktitletext}{#3}%
  \renewcommand{\deckdate}{#4}%
  \renewcommand{\insertshortdeck}{Session #1 \,$\cdot$\, Deck #1.#2}%
  \title{#3}%
  \author{Zhigang Feng}%
  \date{#4}%
}
\newcommand{\makedecktitle}{%
  \begin{frame}[plain,noframenumbering]
    \begin{tikzpicture}[remember picture,overlay]
      \fill[UCRBlue] (current page.north west) rectangle ([yshift=-0.55cm]current page.north east);
      \fill[UCRGold] ([yshift=-0.55cm]current page.north west) rectangle ([yshift=-0.62cm]current page.north east);
      \node[anchor=west, text=white, font=\small\bfseries] at ([xshift=5mm,yshift=-0.275cm]current page.north west)
        {ECON 282E \,$\cdot$\, Foundations of Macroeconomics};
      \node[anchor=east, text=white, font=\small] at ([xshift=-5mm,yshift=-0.275cm]current page.north east)
        {UC Riverside \,$\cdot$\, Fall 2026};
    \end{tikzpicture}
    \vspace*{1.1cm}
    {\usebeamerfont{title}\color{black!55}\large Session \decksession\ \,$\cdot$\, Deck \decksession.\deckid\par}
    \vspace{0.25cm}
    {\usebeamerfont{title}\color{RedTitle}\LARGE\bfseries \decktitletext\par}
    \vspace{0.7cm}
    {\large Zhigang Feng\par}
    \vspace{0.15cm}
    {\color{black!65}\deckdate\par}
    \vfill
    {\footnotesize\itshape\color{black!55}Build it on paper $\;\to\;$ solve it on the machine $\;\to\;$ push it to the frontier with AI\par}
  \end{frame}}

% ---------------------------------------------------------------------------
% Section divider (plain frame, not counted)
\newcommand{\sectiondivider}[2]{%
\begin{frame}[plain,noframenumbering]
\begin{center}
\vfill
{\Large\textbf{#1}\par}
\vspace{0.35cm}
{\normalsize #2\par}
\vfill
\end{center}
\end{frame}
}

% ---------------------------------------------------------------------------
% Boxes
\newtcolorbox{conceptbox}[1][]{colback=blue!3, colframe=RedTitle!55, boxrule=0.35mm,
  arc=1mm, left=2mm, right=2mm, top=1mm, bottom=1mm, #1}
\newtcolorbox{cautionbox}[1][]{colback=red!3, colframe=red!50!black, boxrule=0.35mm,
  arc=1mm, left=2mm, right=2mm, top=1mm, bottom=1mm, #1}
\newtcolorbox{examplebox}[1][]{colback=green!3, colframe=green!45!black, boxrule=0.35mm,
  arc=1mm, left=2mm, right=2mm, top=1mm, bottom=1mm, #1}
\newtcolorbox{takeawaybox}[1][]{colback=VeryLightGray, colframe=RedTitle!55, boxrule=0.35mm,
  arc=1mm, left=2mm, right=2mm, top=1mm, bottom=1mm, #1}
\newtcolorbox{researchbox}[1][]{enhanced, colback=violet!4, colframe=violet!60!black,
  boxrule=0.35mm, arc=1mm, left=2mm, right=2mm, top=1mm, bottom=1mm,
  fonttitle=\bfseries\small, title={Research in the wild}, #1}

\newcommand{\compacttables}{\renewcommand{\arraystretch}{1.15}}
\newcommand{\src}[1]{{\tiny\color{black!55}#1}}

% ---------------------------------------------------------------------------
%% Course map: the recurring "you are here" slide for ECON 282E.
%% Loaded by course_preamble.tex. Pattern follows AI-ECON-2026 lec01_map.tex, but the map
%% shows the whole ten-session course rather than one lecture.
%%
%%   \CourseMap{s}{label}   s = 1..10 highlights session s and prints "you are here: label"
%%                          (e.g. \CourseMap{1}{1.A2}); earlier sessions get a check mark.
%%   \CourseMap{0}{}        full overview, nothing highlighted.
%%
%% Note: TikZ option lists are comma-split before TeX conditionals run, so every \ifnum
%% below sits outside [...] option lists.

\newcommand{\CourseMap}[2]{%
\begin{center}
\begin{tikzpicture}[
    sess/.style={rectangle, rounded corners=2pt, draw=black!45, minimum width=1.34cm,
        minimum height=1.12cm, text width=1.26cm, align=center, inner sep=1pt},
    stage/.style={font=\tiny\bfseries, text=black!70},
]
  \foreach \i/\d/\t/\c in {%
      1/Sep 24/Intro \\ Growth/StageTrunk,
      2/Oct 1/HA $\cdot$ OLG \\ Policy/StageBranch,
      3/Oct 8/Num.\ DP \\ Python/StageMachine,
      4/Oct 15/Numerics \\ I/StageMachine,
      5/Oct 22/Numerics \\ II/StageMachine,
      6/Oct 29/The wall \\ \textbf{Midterm}/StageWall,
      7/Nov 5/ML \\ Deep nets/StageTools,
      8/Nov 12/RL \\ HA $+$ DL/StageTools,
      9/Nov 19/Text \\ LLMs/StageData,
      10/Dec 3/Agents \\ \textbf{Projects}/StageAgents} {
    \node[sess, fill=\c] (n\i) at ({1.46*(\i-1)}, 0)
      {{\scriptsize\bfseries S\i}\\[-1pt]{\tiny\color{black!60}\d}\\[1pt]{\tiny \t}};
    \ifnum\i<#1
      \node[font=\scriptsize, text=green!45!black] at ([xshift=-2.5mm,yshift=-2.2mm]n\i.north east) {$\checkmark$};
    \fi
    \ifnum\i=#1
      \draw[RedTitle, very thick, rounded corners=2pt]
        ([xshift=-0.6mm,yshift=-0.6mm]n\i.south west) rectangle ([xshift=0.6mm,yshift=0.6mm]n\i.north east);
      % keep the label inside the picture at both ends of the row
      \ifnum\i<3
        \node[font=\scriptsize\bfseries, text=RedTitle, anchor=south west] at ([yshift=1.2mm]n\i.north west)
          {$\blacktriangledown$\,you are here: #2};
      \else\ifnum\i>8
        \node[font=\scriptsize\bfseries, text=RedTitle, anchor=south east] at ([yshift=1.2mm]n\i.north east)
          {you are here: #2\,$\blacktriangledown$};
      \else
        \node[font=\scriptsize\bfseries, text=RedTitle, anchor=south] at ([yshift=1.2mm]n\i.north)
          {you are here: #2\,$\blacktriangledown$};
      \fi\fi
    \fi
  }
  % stage band under the sessions
  \foreach \a/\b/\lab/\c in {1/1/trunk/StageTrunk, 2/2/branches/StageBranch,
      3/5/the machine $\cdot$ classical methods/StageMachine, 6/6/the wall/StageWall,
      7/8/new tools/StageTools, 9/9/new data/StageData, 10/10/colleagues/StageAgents} {
    \fill[\c] ([yshift=-1.5mm]n\a.south west) rectangle ([yshift=-4.6mm]n\b.south east);
    \path ([yshift=-1.5mm]n\a.south west) -- ([yshift=-4.6mm]n\b.south east)
      node[midway, stage] {\lab};
  }
  \node[font=\footnotesize\itshape, text=black!70] at ([yshift=-9.5mm]$(n1.south)!0.5!(n10.south)$)
    {Build it on paper $\;\to\;$ solve it on the machine $\;\to\;$ push it to the frontier with AI};
\end{tikzpicture}
\end{center}
}


\graphicspath{{./figures/}}

\lstset{language=Python, basicstyle=\ttfamily\scriptsize, frame=none,
        backgroundcolor=\color{gray!8}, aboveskip=2pt, belowskip=2pt}

\decktitle{4}{B4}{Fast and Accurate: The Same Model, Measured}{Thursday, October 15, 2026}

\begin{document}

\makedecktitle

%=============================================================================
\begin{frame}{Where We Are}
\CourseMap{4}{4.B4}
\vspace{-1mm}
\begin{center}
\small \textbf{Block B:} \textbf{4.B1} approximation and interpolation $\;\cdot\;$ \textbf{4.B2} quadrature $\;\cdot\;$ \textbf{4.B3} discretizing an AR(1) $\;\cdot\;$ \textbf{4.B4} fast and accurate $\;\cdot\;$ \textbf{4.B5} the worked example, end to end.
\end{center}
\end{frame}

%=============================================================================
\begin{frame}{One Model, Six Solvers}
\begin{columns}[T]
\begin{column}{0.56\textwidth}
\small
Quarterly RBC, $n_k=500$, $n_z=7$, tolerance $10^{-6}$, identical in every row.
\medskip
\begin{center}\scriptsize
\begin{tabularx}{\linewidth}{@{}>{\raggedright\arraybackslash}X r r r@{}}
\toprule
 & \textbf{time} & \textbf{speed-up} & \textbf{worst $\log_{10}\mathcal{E}$} \\
\midrule
plain VFI & 3.46\,s & 1.0 & $-1.55$ \\
$+$ Howard ($n_H{=}20$) & 0.25\,s & 14.1 & $-1.55$ \\
$+$ Howard ($n_H{=}50$) & 0.16\,s & 22.1 & $-1.55$ \\
$+$ monotonicity & 1.10\,s & 3.2 & $-1.55$ \\
$+$ continuous choice & 1.11\,s & 3.1 & $-2.77$ \\
\textbf{EGM} & \textbf{0.09\,s} & \textbf{37.7} & $\mathbf{-6.06}$ \\
\bottomrule
\end{tabularx}
\end{center}
\end{column}
\begin{column}{0.41\textwidth}
\begin{takeawaybox}[title=\textbf{Read the two columns separately}]\small
The acceleration rows are \textbf{faster at identical accuracy}. The last two rows are \textbf{more accurate}. Only one row is both.
\end{takeawaybox}
\medskip
\begin{examplebox}\footnotesize
EGM is 38 times faster \emph{and} four and a half decades more accurate --- and it does zero root-finding to get there.
\end{examplebox}
\end{column}
\end{columns}
\end{frame}

%=============================================================================
\begin{frame}{Where the Time Actually Goes}
\begin{columns}[T]
\begin{column}{0.52\textwidth}
\small
The plain solver's bill factorises:
\[
  \underbrace{n_k\,n_z\,n_k}_{\text{one sweep}}\times\underbrace{\frac{\ln\varepsilon}{\ln\beta}}_{\text{patience}}
  \;=\; 1.75\text{M}\times1352 \;=\; \mathbf{2.4\ \text{billion}} .
\]
Two independent things to attack:
\begin{itemize}\footnotesize
  \item \textbf{The sweep count}, which belongs to $\beta$ alone --- 1{,}352 measured against a theoretical 1{,}375, at $\beta=0.99$.
  \item \textbf{The cost per sweep}, which is the maximization over $n_k$ candidates at every one of $n_k n_z$ states.
\end{itemize}
\end{column}
\begin{column}{0.45\textwidth}
\begin{cautionbox}[title=\textbf{This got worse since Session 3}]\small
The same solver needed \textbf{228} sweeps in 3.A4 at $\beta=0.95$. Moving to a quarterly calibration multiplied the patience bill by six, and nothing about the economics changed.
\end{cautionbox}
\medskip
\begin{conceptbox}\footnotesize
Howard attacks the first. Monotonicity and concavity attack the second. \textbf{EGM removes the maximization entirely.}
\end{conceptbox}
\end{column}
\end{columns}
\vspace{1mm}
\src{Cost accounting adapted from QM Lec.~8, ``Recap: Standard VFI'' and ``Computational Bottlenecks''.}
\end{frame}

%=============================================================================
\begin{frame}{Howard: Stop Maximizing So Often}
\begin{columns}[T]
\begin{column}{0.53\textwidth}
\small
The policy settles long before the value does --- 3.A1 measured it, at iteration 11 out of 272. So most sweeps re-derive a policy that has not changed.
\medskip

\textbf{Howard's policy improvement} alternates:
\begin{itemize}\footnotesize
  \item one \emph{improvement} step: the full $\max$, $O(n_k^{2}n_z)$;
  \item then $n_H$ \emph{evaluation} steps, which reuse the stored policy index and cost only $O(n_k n_z)$:
\end{itemize}
\[
  v^{(n+1)}(k_i,z_j)=u\big(c_{ij}\big)+\beta\!\sum_m\!\Pi_{jm}v^{(n)}\!\big(g(k_i,z_j),z_m\big).
\]
No \texttt{argmax} --- just a lookup.
\end{column}
\begin{column}{0.44\textwidth}
\begin{conceptbox}[title=\textbf{Why it converges}]\small
Policy evaluation solves a \emph{linear} fixed point, $v_g=(I-\beta P_g)^{-1}u_g$. Iterating it $n_H$ times is itself a contraction with modulus $\beta$, so you land within $\beta^{\,n_H}$ of $v_g$.
\end{conceptbox}
\medskip
\begin{examplebox}\footnotesize
The improvement step is monotone and there are finitely many grid policies, so the algorithm terminates in finitely many policy updates --- something plain VFI cannot promise.
\end{examplebox}
\end{column}
\end{columns}
\vspace{1mm}
\src{Adapted from QM Lec.~8, ``Howard's Policy Improvement'' block. Step count written $n_H$; the source uses $H$, which this course reserves for the law of motion.}
\end{frame}

%=============================================================================
\begin{frame}{Howard, Measured}
\begin{columns}[T]
\begin{column}{0.52\textwidth}
\begin{center}\footnotesize
\begin{tabularx}{\linewidth}{@{}>{\raggedright\arraybackslash}X r r r@{}}
\toprule
 & \textbf{policy updates} & \textbf{comparisons} & \textbf{time} \\
\midrule
plain VFI & 1{,}352 & 2.37\,bn & 3.46\,s \\
$n_H=10$ & 124 & 221\,m & 0.39\,s \\
$n_H=20$ & 66 & 120\,m & 0.25\,s \\
$n_H=50$ & \textbf{29} & \textbf{56\,m} & \textbf{0.16\,s} \\
\bottomrule
\end{tabularx}
\end{center}
\medskip
\small
From 1{,}352 maximizations to \textbf{29}, at \emph{identical} accuracy --- every row has the same worst Euler error, $10^{-1.55}$.
\end{column}
\begin{column}{0.45\textwidth}
\begin{takeawaybox}[title=\textbf{Diminishing returns}]\small
$n_H:10\to20$ halves the updates. $20\to50$ halves them again but only cuts the time by a third, because the evaluation steps are no longer free.
\end{takeawaybox}
\medskip
\begin{cautionbox}\footnotesize
\textbf{Accuracy did not move.} Howard changes how you reach the fixed point, not which fixed point it is. The discretization error is untouched.
\end{cautionbox}
\end{column}
\end{columns}
\vspace{1mm}
\src{Rule of thumb from the source: $n_H\in[10,50]$. Our measurement supports it.}
\end{frame}

%=============================================================================
\begin{frame}{Monotonicity: Fewer Comparisons, More Wall Clock}
\begin{columns}[T]
\begin{column}{0.52\textwidth}
\small
$g(k,z)$ is increasing in $k$, so the maximizer at $k_{i+1}$ cannot lie below the maximizer at $k_i$. Search upward from wherever the previous state stopped, and the sweep falls from $O(n_k^{2})$ to $O(n_k)$.
\medskip
\begin{center}\footnotesize
\begin{tabularx}{\linewidth}{@{}>{\raggedright\arraybackslash}X r r@{}}
\toprule
Howard $n_H=20$ & \textbf{comparisons} & \textbf{time} \\
\midrule
without monotonicity & 120\,m & \textbf{0.25\,s} \\
with monotonicity & \textbf{64\,m} & 1.10\,s \\
\bottomrule
\end{tabularx}
\end{center}
\medskip
\textbf{Half the comparisons, and four and a half times slower.}
\end{column}
\begin{column}{0.45\textwidth}
\begin{cautionbox}[title=\textbf{Both numbers are true}]\small
The operation count is what the theory predicts. The wall clock is what the machine does --- and a monotone ratchet is a \emph{sequential} loop, while the full search is one vectorised array operation.
\end{cautionbox}
\medskip
\begin{takeawaybox}\footnotesize
\textbf{An algorithmic improvement is not automatically a speed-up.} Report the operation count, which is implementation-independent, \emph{and} the wall clock, which is what you actually wait for. In a compiled language this row goes the other way.
\end{takeawaybox}
\end{column}
\end{columns}
\vspace{1mm}
\src{Algorithm from QM Lec.~8, ``Exploiting Monotonicity''; the wall-clock result is ours, and is not in the source.}
\end{frame}

%=============================================================================
\begin{frame}{Concavity, and the Callback to This Morning}
\begin{columns}[T]
\begin{column}{0.53\textwidth}
\small
Concavity of $u$ and of the interpolated $\hat v$ makes the inner objective \textbf{unimodal}, which is exactly the licence golden section needs --- and 4.B1 established that linear interpolation is what supplies it.
\medskip

So the inner maximization can be replaced by 4.A2's bracket-shrinking, at $O(\log(1/\varepsilon))$ evaluations instead of $O(n_k)$ comparisons, \emph{and} the answer is no longer confined to a grid point.
\medskip

Measured: the Euler error falls from $10^{-1.55}$ to $\mathbf{10^{-2.77}}$, and the objective-evaluation count falls from 2.4 billion to \textbf{18.9 million}.
\end{column}
\begin{column}{0.44\textwidth}
\begin{conceptbox}[title=\textbf{Two gains, one change}]\small
Monotonicity and Howard made the \emph{same} answer cheaper. Continuous choice makes a \emph{different}, better answer --- it moves the accuracy, not just the clock.
\end{conceptbox}
\medskip
\begin{examplebox}\footnotesize
Combine them: use monotonicity to bracket, then golden section inside the bracket, then Howard around the whole thing.
\end{examplebox}
\end{column}
\end{columns}
\vspace{1mm}
\src{Adapted from QM Lec.~8, ``Exploiting Concavity'' and ``Golden Section Search in VFI''. The source's ``Recall Lecture 4'' does not apply here --- golden section was deliberately withheld until 4.A2.}
\end{frame}

%=============================================================================
\begin{frame}{The Endogenous Grid Method: Run the Euler Equation Backwards}
\begin{columns}[T]
\begin{column}{0.53\textwidth}
\small
Time iteration is expensive because $k'$ appears on \emph{both} sides of
\[
  u'\big(y(k,z)-k'\big)=\beta\,\mathbb{E}\big[u'(c(k',z'))R(k',z')\big],
\]
so every state needs a root-find, every sweep.
\medskip

\textbf{Carroll's insight (2006):} treat the grid as \emph{tomorrow's} assets. Fix $k'$. Then the right-hand side is computable immediately, and the left-hand side inverts:
\[
  c=(u')^{-1}\Big(\beta\,\mathbb{E}\big[u'(c(k',z'))R(k',z')\big]\Big).
\]
Now you know $c$ and $k'$ --- and today's state follows.
\end{column}
\begin{column}{0.44\textwidth}
\begin{conceptbox}[title=\textbf{The trade}]\small
The unknown moves from \emph{inside} a non-linear equation to the \emph{outside} of an explicit formula. For CRRA, $(u')^{-1}(x)=x^{-1/\sigma}$ --- one exponentiation.
\end{conceptbox}
\medskip
\begin{examplebox}\footnotesize
The grid you end up with is \textbf{endogenous}: the $k$ values that generate your chosen $k'$ are wherever the algebra puts them. One interpolation brings the policy back to the grid you wanted.
\end{examplebox}
\end{column}
\end{columns}
\vspace{1mm}
\src{Carroll, ``The method of endogenous gridpoints for solving dynamic stochastic optimization problems'', \emph{Economics Letters} (2006).}
\end{frame}

%=============================================================================
\begin{frame}{Getting EGM Right: Put the Grid on Cash on Hand}
\begin{columns}[T]
\begin{column}{0.53\textwidth}
\small
Almost every presentation of EGM says ``no root-finding required'' --- and then, with Cobb--Douglas production, has to solve
\[
  f(\tilde k)+(1-\delta)\tilde k-c-k'=0
\]
for today's capital $\tilde k$. That is a root-find, at every node, every sweep. The headline claim is retracted three slides after it is made.
\medskip

\textbf{The fix is to choose the right state.} Let
\[
  m=z k^{\alpha}+(1-\delta)k \quad\text{(cash on hand)} .
\]
Then the budget constraint is simply $m=c+k'$, so
\[
  \tilde m = c + k' ,
\]
with no equation to solve at all.
\end{column}
\begin{column}{0.44\textwidth}
\begin{takeawaybox}[title=\textbf{Why this is the real method}]\small
Carroll's formulation is in cash on hand for exactly this reason. Keeping the grid in $k$ reintroduces the root-find the method exists to remove.
\end{takeawaybox}
\medskip
\begin{cautionbox}\footnotesize
Measured consequence: our cash-on-hand EGM does \textbf{zero} root solves and runs in 0.09\,s. An implementation that keeps the grid in $k$ does 3{,}500 root solves per sweep and loses the speed race it was supposed to win.
\end{cautionbox}
\end{column}
\end{columns}
\end{frame}

%=============================================================================
\begin{frame}{EGM in Four Steps}
\begin{columns}[T]
\begin{column}{0.54\textwidth}
\small
For each $z_j$ and each end-of-period $k'$ on the exogenous grid:
\medskip
\begin{enumerate}\footnotesize
  \item \textbf{Next period's cash} for every $z'$:
        $m'=z'(k')^{\alpha}+(1-\delta)k'$, and $c'=\hat c(m',z')$ by interpolation.
  \item \textbf{The expectation}:
        $\;\mathcal{R}=\beta\sum_{m}\Pi_{jm}\,u'(c')\,R(k',z_m)$.
  \item \textbf{Invert}: $\;c=(u')^{-1}(\mathcal{R})$. \;For log utility, $c=1/\mathcal{R}$.
  \item \textbf{Today's state}: $\;\tilde m=c+k'$. Interpolate $(\tilde m,c)$ back onto the $m$-grid.
\end{enumerate}
\medskip
\footnotesize
Where $m$ falls below the smallest $\tilde m$, the constraint binds and $c=m-\underline{k}$.
\end{column}
\begin{column}{0.43\textwidth}
\begin{conceptbox}[title=\textbf{Count the solves}]\small
Steps 1--4 contain one interpolation, one matrix product, one exponentiation and one addition. \textbf{No maximization and no root-finding.}
\end{conceptbox}
\medskip
\begin{examplebox}\footnotesize
That is the whole method. It is shorter than the VFI loop it replaces, and it converged here in 420 iterations and 0.09 seconds.
\end{examplebox}
\end{column}
\end{columns}
\vspace{1mm}
\src{Steps adapted from QM Lec.~8's EGM block, restated on cash on hand.}
\end{frame}

%=============================================================================
\begin{frame}{When EGM Does Not Apply}
\begin{columns}[T]
\begin{column}{0.52\textwidth}
\small
The method needs four things, and each is a real restriction:
\begin{enumerate}\footnotesize
  \item $u'$ must be \textbf{invertible} in closed form --- CRRA, CARA, quadratic. Not satiation.
  \item The budget must solve for today's state \textbf{explicitly} given $(c,k')$ --- which is what cash on hand buys.
  \item \textbf{One} continuous control. Two, and part of the problem goes back to root-finding.
  \item The endogenous grid must come out \textbf{ordered}, which needs a monotone policy.
\end{enumerate}
\end{column}
\begin{column}{0.45\textwidth}
\begin{cautionbox}[title=\textbf{What breaks item 4}]\small
Discrete choices --- to work or not, to default or not --- make the policy non-monotone, the endogenous grid folds over itself, and a plain interpolation picks the wrong branch.
\end{cautionbox}
\medskip
\begin{examplebox}\footnotesize
The repair is an \textbf{upper-envelope} step: Fella (2014), and Iskhakov, J\o rgensen, Rust \& Schjerning (2017) for the discrete-continuous case. Both are outside this course, and both are what you would read next.
\end{examplebox}
\end{column}
\end{columns}
\vspace{1mm}
\src{Conditions adapted from QM Lec.~8, ``When It Works and When It Doesn't''; the DC-EGM reference is added.}
\end{frame}

%=============================================================================
\begin{frame}{The Accuracy, Where the Source Has No Figure}
\begin{columns}[c]
\begin{column}{0.53\textwidth}
\centering
\begin{tikzpicture}
\begin{axis}[
    width=7.9cm, height=5.5cm, xmin=14, xmax=42.8, ymin=-7.4, ymax=-1.0,
    xlabel={\scriptsize capital $k$ (off the solution grid)},
    ylabel={\scriptsize $\log_{10}\mathcal{E}$},
    tick label style={font=\scriptsize},
    axis x line*=bottom, axis y line*=left,
    legend style={font=\tiny, draw=none, fill=none, at={(0.98,0.55)}, anchor=east}]
  \addplot[thin, RedTitle] coordinates {(14.2878,-1.5946) (14.5697,-1.9789) (14.8517,-2.0228) (15.1336,-2.0414) (15.4155,-2.1460) (15.6974,-2.1867) (15.9793,-1.6159) (16.2612,-1.6116) (16.5432,-2.0777) (16.8251,-2.0978) (17.1070,-2.1183) (17.3889,-2.2564) (17.6708,-2.2845) (17.9527,-1.6239) (18.2347,-2.0945) (18.5166,-2.1685) (18.7985,-2.1912) (19.0804,-2.3495) (19.3623,-2.3824) (19.6442,-1.6302) (19.9262,-2.2306) (20.2081,-2.2556) (20.4900,-2.2815) (20.7719,-2.3279) (21.0538,-2.5151) (21.3357,-1.6325) (21.6177,-2.3347) (21.8996,-2.3643) (22.1815,-2.3954) (22.4634,-2.4282) (22.7453,-2.7091) (23.0272,-3.2150) (23.3092,-2.4687) (23.5911,-2.5065) (23.8730,-2.5471) (24.1549,-2.8029) (24.4368,-1.8310) (24.7187,-2.4026) (25.0007,-2.6567) (25.2826,-2.7120) (25.5645,-2.7741) (25.8464,-4.3714) (26.1283,-1.9144) (26.4102,-2.6622) (26.6922,-2.9787) (26.9741,-3.0943) (27.2560,-3.2491) (27.5379,-2.9939) (27.8198,-3.3895) (28.1017,-3.7668) (28.3837,-4.2085) (28.6656,-3.5360) (28.9475,-1.6424) (29.2294,-3.0589) (29.5113,-3.1594) (29.7932,-3.0390) (30.0752,-2.9462) (30.3571,-2.5833) (30.6390,-2.4285) (30.9209,-2.8129) (31.2028,-2.7572) (31.4847,-2.7085) (31.7667,-2.6654) (32.0486,-1.9581) (32.3305,-2.9643) (32.6124,-2.5908) (32.8943,-2.5586) (33.1762,-2.5290) (33.4582,-2.1631) (33.7401,-2.7158) (34.0220,-2.4733) (34.3039,-2.4497) (34.5858,-2.4275) (34.8677,-2.4067) (35.1497,-1.7826) (35.4316,-2.3830) (35.7135,-2.3646) (35.9954,-2.3471) (36.2773,-2.3305) (36.5592,-1.7857) (36.8412,-2.4310) (37.1231,-2.2951) (37.4050,-2.2808) (37.6869,-2.2671) (37.9688,-1.8190) (38.2507,-3.0533) (38.5327,-2.2367) (38.8146,-2.2247) (39.0965,-2.2131) (39.3784,-1.8536) (39.6603,-1.8800) (39.9422,-2.1864) (40.2242,-2.1761) (40.5061,-2.1662) (40.7880,-1.8897) (41.0699,-1.8949) (41.3518,-2.1425) (41.6337,-2.1336) (41.9157,-2.1249) (42.1976,-1.9276)};
  \addlegendentry{plain VFI}
  \addplot[thin, orange!80!black] coordinates {(14.2878,-6.2038) (14.5697,-6.1200) (14.8517,-6.1371) (15.1336,-6.2497) (15.4155,-6.3127) (15.6974,-6.3973) (15.9793,-6.5100) (16.2612,-6.4013) (16.5432,-6.4337) (16.8251,-6.4324) (17.1070,-6.3935) (17.3889,-6.4713) (17.6708,-6.7521) (17.9527,-6.7210) (18.2347,-6.5470) (18.5166,-6.4041) (18.7985,-6.3830) (19.0804,-6.4538) (19.3623,-6.6622) (19.6442,-6.8107) (19.9262,-6.5215) (20.2081,-6.4214) (20.4900,-6.4171) (20.7719,-6.4962) (21.0538,-6.4577) (21.3357,-6.4792) (21.6177,-6.5956) (21.8996,-6.6311) (22.1815,-6.6510) (22.4634,-6.5568) (22.7453,-6.5639) (23.0272,-6.6687) (23.3092,-6.9672) (23.5911,-7.2905) (23.8730,-7.3435) (24.1549,-6.9134) (24.4368,-6.8006) (24.7187,-6.8282) (25.0007,-6.6715) (25.2826,-6.6363) (25.5645,-6.6861) (25.8464,-6.8488) (26.1283,-7.0348) (26.4102,-6.8818) (26.6922,-6.7805) (26.9741,-6.7962) (27.2560,-6.9243) (27.5379,-6.8672) (27.8198,-6.9306) (28.1017,-6.9438) (28.3837,-7.0670) (28.6656,-7.1125) (28.9475,-6.8308) (29.2294,-6.7170) (29.5113,-6.7162) (29.7932,-6.7641) (30.0752,-6.7288) (30.3571,-6.7665) (30.6390,-6.8942) (30.9209,-7.2424) (31.2028,-7.7871) (31.4847,-7.4024) (31.7667,-7.0801) (32.0486,-7.0032) (32.3305,-6.9329) (32.6124,-6.8433) (32.8943,-6.8402) (33.1762,-6.9167) (33.4582,-7.0911) (33.7401,-7.1543) (34.0220,-7.4398) (34.3039,-7.6645) (34.5858,-7.4084) (34.8677,-7.0683) (35.1497,-6.9232) (35.4316,-6.8822) (35.7135,-6.9143) (35.9954,-6.9307) (36.2773,-6.9367) (36.5592,-7.0264) (36.8412,-7.2829) (37.1231,-8.0002) (37.4050,-7.2576) (37.6869,-7.0654) (37.9688,-7.0098) (38.2507,-7.0388) (38.5327,-6.9687) (38.8146,-6.9767) (39.0965,-7.0626) (39.3784,-7.2919) (39.6603,-7.5443) (39.9422,-7.5887) (40.2242,-7.5182) (40.5061,-7.3361) (40.7880,-7.1841) (41.0699,-7.0761) (41.3518,-7.0640) (41.6337,-7.1318) (41.9157,-7.2437) (42.1976,-7.3220)};
  \addlegendentry{EGM}
  \draw[black!45, dashed] (axis cs:14,-3) -- (axis cs:42.8,-3);
  \node[font=\tiny, text=black!55, anchor=south east] at (axis cs:42.5,-2.95) {``acceptable'' $=10^{-3}$};
\end{axis}
\end{tikzpicture}
\end{column}
\begin{column}{0.44\textwidth}
\small
Unit-free Euler residuals at 400 capital levels \emph{between} the solution grid points, at the middle productivity state.
\medskip
\begin{center}\footnotesize
\begin{tabularx}{\linewidth}{@{}>{\raggedright\arraybackslash}X r r@{}}
\toprule
 & \textbf{mean} & \textbf{worst} \\
\midrule
plain VFI & $-2.18$ & $-1.55$ \\
EGM & $\mathbf{-6.75}$ & $\mathbf{-6.06}$ \\
\bottomrule
\end{tabularx}
\end{center}
\medskip
\begin{takeawaybox}\footnotesize
\textbf{Four and a half decades}, on the same grid, the same chain and the same tolerance. VFI's worst case is a 2.8\% consumption mistake; EGM's is one part in a million.
\end{takeawaybox}
\end{column}
\end{columns}
\vspace{1mm}
\src{QM Lec.~8's ``Euler Errors: Visualization'' frame describes this plot in words and contains no figure. This one is drawn from our run.}
\end{frame}

%=============================================================================
\begin{frame}{Reporting: What Goes in the Paper}
\begin{columns}[T]
\begin{column}{0.51\textwidth}
\small
\begin{enumerate}\footnotesize
  \item The method, and the implementation choices that matter.
  \item The grid: points, spacing, bounds --- and the chain, with its \emph{implied} $\rho$ and $\sigma_z$ (4.B3).
  \item The convergence tolerance.
  \item \textbf{Maximum and mean Euler errors on an evaluation grid that is not the solution grid.}
  \item A grid-density study, and whether the observed rate matches theory.
  \item Computation time and hardware.
\end{enumerate}
\end{column}
\begin{column}{0.46\textwidth}
\begin{cautionbox}[title=\textbf{Item 4 is the one with teeth}]\small
Evaluating residuals \emph{on} the solution grid is the numerical analogue of reporting training accuracy. The solution was constructed to satisfy the Euler equation there.
\end{cautionbox}
\medskip
\begin{takeawaybox}\footnotesize
And item 5 has to name which object the rate applies to. 3.A2 measured \textbf{1.89 for the value function and 0.93 for the policy} --- quoting a single rate without saying which is how the two get confused.
\end{takeawaybox}
\end{column}
\end{columns}
\vspace{1mm}
\src{Checklist adapted from QM Lec.~8, ``Reporting Standards''; Den Haan--Marcet is the simulation-based complement, and needs a stochastic model to be non-degenerate.}
\end{frame}

%=============================================================================
\begin{frame}{The Benchmark Table, Recomputed}
\renewcommand{\arraystretch}{1.15}
\begin{center}\scriptsize
\begin{tabularx}{\linewidth}{@{}>{\raggedright\arraybackslash}p{3.4cm} r r r >{\raggedright\arraybackslash}X@{}}
\toprule
\textbf{Method} & \textbf{time} & \textbf{speed-up} & \textbf{worst $\log_{10}\mathcal{E}$} & \textbf{what changed} \\
\midrule
Plain VFI, grid search & 3.46\,s & 1.0 & $-1.55$ & --- \\
$+$ Howard, $n_H=20$ & 0.25\,s & 14.1 & $-1.55$ & fewer maximizations \\
$+$ Howard, $n_H=50$ & 0.16\,s & 22.1 & $-1.55$ & fewer still \\
$+$ monotonicity & 1.10\,s & 3.2 & $-1.55$ & fewer comparisons, more overhead \\
$+$ continuous choice & 1.11\,s & 3.1 & $-2.77$ & the policy leaves the grid \\
\textbf{EGM, cash on hand} & \textbf{0.09\,s} & \textbf{37.7} & $\mathbf{-6.06}$ & no max, no root-find \\
\bottomrule
\end{tabularx}
\end{center}
\vspace{1mm}
\begin{columns}[T]
\begin{column}{0.49\textwidth}
\begin{cautionbox}\footnotesize
Published versions of this table are often attributed to Aruoba, Fern\'andez-Villaverde \& Rubio-Ram\'irez (2006). That paper compares VFI, perturbation, Chebyshev projection, finite elements and PEA --- it has \textbf{no EGM row and no time-iteration row}.
\end{cautionbox}
\end{column}
\begin{column}{0.48\textwidth}
\begin{takeawaybox}\footnotesize
So this one is ours, on one machine, at one calibration, with the script in the repository. \textbf{Cite a benchmark table to the run that produced it.}
\end{takeawaybox}
\end{column}
\end{columns}
\end{frame}

%=============================================================================
\begin{frame}{Takeaways}
\begin{takeawaybox}
\begin{itemize}\small
  \item The bill factorises into \textbf{sweeps} (which belong to $\beta$) and \textbf{cost per sweep} (which belongs to the maximization). They are attacked by different tools.
  \item \textbf{Howard} cut 1{,}352 maximizations to 29 at identical accuracy --- a $22\times$ speed-up that does not move the answer at all.
  \item \textbf{Monotonicity halved the comparisons and quadrupled the wall clock}, because a ratchet is sequential and a full search vectorises. Report both numbers; in a compiled language the sign flips.
  \item Continuous choice is the first change that buys \emph{accuracy} rather than speed: $10^{-1.55}\to10^{-2.77}$.
  \item \textbf{EGM wins on both}: 38 times faster and four and a half decades more accurate, with zero root-finding --- but only if the grid is on \textbf{cash on hand}. Keep it in $k$ and the root-find comes back.
  \item Report residuals \textbf{off} the solution grid, and cite a benchmark table to the run that produced it.
\end{itemize}
\end{takeawaybox}
\end{frame}

%=============================================================================
\begin{frame}{Bridge: What You Owe, and What Comes Next}
\begin{columns}[T]
\begin{column}{0.51\textwidth}
\small
\textbf{Session 4 in one sentence.} Every crude choice of Session 3 had a repair; we built each one and measured what it bought --- and the largest single gain came from \emph{reformulating the state}, not from a better algorithm.
\medskip
\begin{itemize}\footnotesize
  \item \textbf{4.B5} puts the whole toolkit into one program: the stochastic model solved again, line by line, against the crude solver.
  \item \textbf{Lab L04} runs all of it: the root-finder bake-off, golden section, the interpolation shape tests, quadrature, Tauchen against Rouwenhorst, and the accelerated RBC.
  \item \textbf{Homework 1, Part B} is out today, due October 22 --- and Part A is due the same day.
\end{itemize}
\end{column}
\begin{column}{0.46\textwidth}
\begin{conceptbox}[title=\textbf{Next week (S5)}]\small
Numerical Methods II. Everything so far has put values on a grid. \textbf{Perturbation} throws the grid away and expands around the steady state; \textbf{projection} replaces $N$ separate values with a handful of basis coefficients.\\[1mm]
Then we solve Aiyagari end to end --- and meet the wall that Session 6 is about.
\end{conceptbox}
\vspace{-1mm}
\begin{examplebox}\footnotesize
And two threads from 4.A3 are left open on purpose: \textbf{CG and SGD} return in \textbf{S7}, as the optimizers that actually train a network, and the \textbf{Metropolis--Hastings} rule returns in \textbf{S9}, where the object being sampled is a posterior rather than a maximum.
\end{examplebox}
\end{column}
\end{columns}
\end{frame}

\end{document}
